\section{Injective moments and a dichotomy for one Hadamard matrix}
\label{sec:pair}

For a matrix $H$ and a row permutation $\pi$, write $H^\pi(i,k)=H(\pi(i),k)$.
The following dichotomy is the single-matrix step in the proof of
Theorem~\ref{thm:fourier}.

\begin{corollary}[Dichotomy]
\label{cor:dichotomy}
For every $6\times6$ Hadamard matrix $H$, at least one of the following
conditions holds:
\begin{enumerate}[label=\textup{(\Alph*)},ref=\Alph*]
  \item[(M)] $g_{H^\pi}(\alpha)=0$ for every row permutation $\pi$;
  \item[(C)] the entrywise square $H^{\circ2}$ is a Hadamard matrix.
\end{enumerate}
\end{corollary}

Condition (M) is the Fourier condition in
Matolcsi, Ruzsa, and Weiner's conjecture
\cite[Conjecture~2.3]{MatolcsiRuzsaWeiner2013}; Section~\ref{sec:cubic}
will classify condition (C). We prove the dichotomy from an exact identity
between averages of nonnegative quantities, using only orthogonality and
unimodularity of one matrix.

\subsection{Injective moments}

Fix a $6\times6$ Hadamard matrix $H$.  An \emph{exponent matrix} is an
integer matrix $A=(A_{uv})$ whose row sums and column sums are all zero.
For $A$ of size $p\times q$, with $p,q\leq6$, define
\begin{equation}
  m_H(A)=
  \E\,\Ree\!\left[
    \prod_{u=0}^{p-1}\prod_{v=0}^{q-1}
      K\bigl(\rho(u),\kappa(v)\bigr)^{A_{uv}}
  \right].
  \label{eq:injective-moment}
\end{equation}
Here $K$ is $H$ or $H^{\mathsf T}$ with probability $1/2$ each,
$\rho:\{0,\ldots,p-1\}\to\{0,\ldots,5\}$ and
$\kappa:\{0,\ldots,q-1\}\to\{0,\ldots,5\}$ are uniformly chosen
injections, and these three choices are independent.  Negative exponents
are well defined because every matrix entry has modulus one.  We set
$m_H(\varnothing)=1$ for the empty exponent matrix.  In particular,
\begin{equation}
  |m_H(A)|\leq1.
  \label{eq:injective-bound}
\end{equation}

Permuting the rows or columns of $A$ leaves this expectation unchanged.
Replacing $A$ by $-A$ conjugates the monomial and therefore leaves its real
part unchanged.  Also $m_H(A^{\mathsf T})=m_H(A)$: transposition exchanges
the two injections and interchanges the equally weighted choices of $K$.
Finally, a zero row or column can be deleted.  Indeed, its factor is one,
and deleting a coordinate of a uniformly chosen injection leaves a
uniformly chosen injection on the remaining coordinates.  These facts
justify all the identifications of exponent matrices used below.

For a balanced integer row $a$, its order is
$\ord(a)=\sum_{j:a_j>0}a_j$, consistently with the notation for charges.
The following relation eliminates any row of order one.

\begin{lemma}[Contraction of an order-one row]
\label{lem:contraction}
Let $A$ be a balanced exponent matrix with $p$ rows and $q\leq6$ columns,
and suppose row $i$ has order one.  For $\ell\ne i$, let
$A_{i\to\ell}$ be obtained by adding row $i$ to row $\ell$ and then
deleting row $i$.  If $p\leq6$, then
\begin{equation}
  (7-p)m_H(A)
    =-\sum_{\substack{0\leq\ell<p\\\ell\ne i}}
        m_H(A_{i\to\ell}).
  \label{eq:injective-contraction}
\end{equation}
If $p=7$, the valid relation is instead
\begin{equation}
  \sum_{\substack{0\leq\ell<7\\\ell\ne i}}
      m_H(A_{i\to\ell})=0.
  \label{eq:seven-row-contraction}
\end{equation}
Every moment in \eqref{eq:seven-row-contraction} has at most six rows;
no moment with seven injectively sampled rows is defined or used.
The analogous relations hold for columns by transposition.
\end{lemma}

\begin{proof}
Since row $i$ is balanced and has order one, it has a $+1$ in one column
$v_+$, a $-1$ in another column $v_-$, and zeros elsewhere.  Fix $K$, the
column injection $\kappa$, and an injective assignment of the other
$p-1$ rows.  Write $M$ for the product of the factors contributed by
those rows.  The columns $\kappa(v_+)$ and $\kappa(v_-)$ are distinct,
so column orthogonality gives
\begin{equation*}
  M\sum_{t=0}^{5}
     K\bigl(t,\kappa(v_+)\bigr)
     \overline{K\bigl(t,\kappa(v_-)\bigr)}=0.
\end{equation*}
Both choices of $K$ are Hadamard matrices: for a square matrix, column
orthogonality also gives row orthogonality.

When $p\leq6$, there are $7-p$ row labels not yet assigned.  The sum over
those labels is $7-p$ times the conditional expectation obtained by
assigning row $i$ uniformly to an unused label.  Averaging the already
fixed assignments produces $(7-p)m_H(A)$.  For each occupied label,
say that of row $\ell$, the corresponding monomial adds the exponents
of row $i$ to those of row $\ell$.  Its average is
$m_H(A_{i\to\ell})$.  Taking real parts and averaging also over $K$
and $\kappa$ proves \eqref{eq:injective-contraction}.  If a contraction
creates a zero row or column, deletion is justified by the marginal
argument preceding the lemma.

When $p=7$, the six other rows occupy all six available labels.  There
is no unused-label contribution.  The same orthogonality equation is
therefore precisely the sum of the six contractions in
\eqref{eq:seven-row-contraction}.
\end{proof}

Repeatedly applying Lemma~\ref{lem:contraction} to rows or columns of
order one expresses any moment with both dimensions at most six as a
rational linear combination of the empty moment and moments whose
nonzero rows and columns all have order at least two.  Each contraction
reduces the number of rows or columns, so this procedure terminates.
Choosing representatives under the symmetries above introduces no
additional mathematical relation.

\subsection{From products of characters to injective moments}

For a permutation $\pi$ of $\{0,\ldots,5\}$, write
$K^\pi(i,k)=K(\pi(i),k)$.  We use the notation
\begin{equation}
  \bigl\langle F(\widetilde H)\bigr\rangle_H
   =\frac{1}{2\cdot6!}
     \sum_{K\in(H,H^{\mathsf T})}\ \sum_{\pi\in S_6} F(K^\pi)
  \label{eq:pair-average}
\end{equation}
for the average over the two orientations and all row permutations.
The ordered pair $(H,H^{\mathsf T})$ in this formula is counted with
multiplicity even when its entries coincide.  Column permutations do
not change $g_{\widetilde H}$ and need not be included.

Let $(6)_k=6\cdot5\cdots(6-k+1)$, with $(6)_0=1$.

\begin{lemma}[Equality-pattern expansion]
\label{lem:partitions}
Let $a_1,\ldots,a_n\in\Z^6$ be balanced charges satisfying
$\sum_{j=1}^n a_j=0$, and let $A$ be the matrix with these columns.
For a set partition $\mathcal S$ of $\{1,\ldots,n\}$, form
$A_{\mathcal S}$ by replacing the columns in each block $S\in\mathcal S$
by their sum $\sum_{j\in S}a_j$.  Then
\begin{equation}
  \left\langle
    \Ree\prod_{j=1}^n g_{\widetilde H}(a_j)
  \right\rangle_H
  =\sum_{\substack{\mathcal S\text{ a set partition of }\{1,\ldots,n\}\\
                    |\mathcal S|\leq6}}
       \frac{(6)_{|\mathcal S|}}{6^n}\,
       m_H(A_{\mathcal S}).
  \label{eq:injective-partitions}
\end{equation}
The ordering of the blocks of $\mathcal S$ is immaterial.
\end{lemma}

\begin{proof}
Expanding the product of characters samples $n$ columns independently
and uniformly from the six columns of $\widetilde H$.  Group the
summands by their equality pattern: two positions belong to the same
block precisely when their column labels coincide.  For a specified
partition with $k$ blocks, the distinct block labels have $(6)_k$
possible assignments, each with probability $6^{-n}$.  Conditional on
this pattern, the block labels form a uniform injection.  Multiplying
the factors with the same column label adds their exponents, giving
$A_{\mathcal S}$.  Averaging over the row permutation, the orientation,
and the real part gives exactly its injective moment.  Partitions with
more than six blocks have no assignments and contribute nothing.
\end{proof}

The hypotheses on the charges ensure that every matrix in
\eqref{eq:injective-partitions} has both row and column sums zero.
Moreover, $\overline{g_H(a)}=g_H(-a)$, so absolute squares of products
are covered by the lemma.  A zero charge may be omitted before applying
the expansion, since its factor is $g_H(0)=1$; equivalently, integrating
out its independently sampled column does not change the distribution
of the remaining columns.

For example, if $a$ is a balanced charge and $[a,-a]$ denotes the
two-column exponent matrix, the two possible equality patterns give
\[
 \left\langle |g_{\widetilde H}(a)|^2\right\rangle_H
 =\frac16+\frac56m_H([a,-a]).
\]
The first term comes from equal sampled columns, where all exponents
cancel; the second comes from two distinct columns. The certificate
below uses the same expansion for products with up to four factors.

\subsection{An exact identity and its consequence}

The auxiliary expression below is chosen so that orthogonality yields a
sum of nonnegative terms equal to zero. The resulting vanishing products
will force the dichotomy.

Retain $\alpha=(1,1,1,-1,-1,-1)$ and put
\begin{align*}
  w&=(-3,-2,0,1,2,2),\\
  a_1&=(-1,-1,0,0,1,1),&
  a_2&=(-2,0,0,0,1,1),\\
  a_3&=(-1,-1,0,1,0,1),&
  a_4&=(-1,-1,0,1,1,0).
\end{align*}
Define
\begin{equation}
  L_H=-g_H(w)+\sum_{j=1}^4 c_j g_H(a_j)g_H(w-a_j),
  \qquad (c_1,c_2,c_3,c_4)=(2,1,1,1),
  \label{eq:pair-L}
\end{equation}
and the two charges
\begin{equation*}
  q_{\mathrm{cross}}=(2,0,0,-2,0,0),
  \qquad q_{\mathrm{same}}=(2,-2,0,0,0,0).
\end{equation*}
Every term defining $L_H$ has total charge $w$.
The nonzero entries of $q_{\mathrm{cross}}$ lie in opposite sign triples
of $\alpha$, whereas those of $q_{\mathrm{same}}$ lie in the same triple.
These two choices will supply the squared-row inner products in the
proof of the dichotomy.

\begin{proposition}[Single-matrix certificate]
\label{prop:pair-certificate}
For every $6\times6$ Hadamard matrix $H$ and for each
$q\in\{q_{\mathrm{cross}},q_{\mathrm{same}}\}$,
\begin{equation}
  3\bigl\langle |L_{\widetilde H}|^2\bigr\rangle_H
  +\left\langle
       |g_{\widetilde H}(\alpha)g_{\widetilde H}(q)|^2
    \right\rangle_H=0.
  \label{eq:pair-certificate}
\end{equation}
\end{proposition}

\begin{proof}
Expand the two absolute squares and apply
Lemma~\ref{lem:partitions} to each resulting product.  Every product has
total charge zero, at most four factors, and sum of the orders of its
charges equal to ten.  Consequently the resulting expression is a
rational linear combination of injective moments.  Apply
Lemma~\ref{lem:contraction} until all remaining nonzero row and column
orders are at least two, and choose representatives under row and
column permutations, negation, and transposition.

The relations used to reduce this expression further are
\eqref{eq:seven-row-contraction}, generated from exponent matrices
with seven rows, two through four columns, positive row orders,
column orders at least two, and total order at most eleven, with a
chosen row of order one.  The total order of an exponent matrix is the
sum of its row orders, equivalently the sum of its column orders.
The right sides of these relations are reduced by the same contraction
procedure.  Thus all relations used are consequences of
Lemma~\ref{lem:contraction} and the stated symmetries.

Section~\ref{sec:verification} verifies in exact rational arithmetic
that substitution of these relations makes every coefficient of the
expanded left side of \eqref{eq:pair-certificate} zero, for each of the
two specified charges $q$.  In particular, its elimination rules hold
for every moment vector satisfying those relations, whether or not
there are further relations between the moments.  The complete finite
computation, including the two coefficient comparisons, is included in
Appendix~\ref{app:code}.  The resulting identity therefore holds for
the moment vector of every Hadamard matrix $H$, which proves the
proposition.
\end{proof}

\begin{proof}[Proof of Corollary~\ref{cor:dichotomy}]
The average in \eqref{eq:pair-certificate} is finite, and each choice
of orientation and row permutation has positive weight.  Every
summand is nonnegative.  Proposition~\ref{prop:pair-certificate}
therefore implies, separately for every ordering of the rows of $H$,
\begin{equation*}
  g_{H^\pi}(\alpha)g_{H^\pi}(q_{\mathrm{cross}})=0,
  \qquad
  g_{H^\pi}(\alpha)g_{H^\pi}(q_{\mathrm{same}})=0.
\end{equation*}
Suppose (M) fails.  The ordering witnessing its failure divides the
row indices into two triples $I_+$ and $I_-$ on which the charge is
$+1$ and $-1$, respectively.  Permutations within either triple do
not change this nonzero character.  By placing any chosen
$i\in I_+$ and $j\in I_-$ in positions $0$ and $3$, the first product
identity gives
$g_H(2(e_i-e_j))=0$.  By placing any two distinct indices of $I_+$ in
positions $0$ and $1$, the second product identity gives the same
vanishing for that pair.  Interchanging the two triples conjugates
the nonzero character, since $g_H(-a)=\overline{g_H(a)}$.
The same argument therefore treats pairs within $I_-$ and all
remaining ordered cross pairs.  Hence
\begin{equation*}
  g_H\bigl(2(e_i-e_j)\bigr)=0
  \qquad\text{for every }i\ne j.
\end{equation*}
For distinct rows $i,j$, this says
\begin{equation*}
  \sum_{k=0}^5 H(i,k)^2\overline{H(j,k)^2}
      =6g_H\bigl(2(e_i-e_j)\bigr)=0.
\end{equation*}
Every row of $H^{\circ2}$ has squared norm six.  Thus
$H^{\circ2}(H^{\circ2})^*=6I$, which for a square matrix also implies
$(H^{\circ2})^*H^{\circ2}=6I$.  Its entries are unimodular, so (C)
holds.
\end{proof}
